\chapter{The LHC and CMS Detector}
\label{ch:detector}

This chapter describes the experimental apparatus used to collect the data analyzed in this thesis. The analysis uses proton-proton collision data recorded during LHC Run~2 (2015--2018) at $\sqrt{s} = 13\TeV$ and Run~3 (2022--2023) at $\sqrt{s} = 13.6\TeV$. The chapter begins with a description of the Large Hadron Collider (LHC), followed by the Compact Muon Solenoid (CMS) detector, the particle-flow event reconstruction algorithm, and the general principles of Monte Carlo simulation.

\section{The Large Hadron Collider}
\label{sec:lhc}

The Large Hadron Collider (LHC)~\cite{Evans:2008zzb} is a circular hadron collider located at CERN (European Organization for Nuclear Research) near Geneva, Switzerland. It occupies a 26.7~km circumference tunnel, originally constructed for the Large Electron-Positron Collider (LEP), at depths ranging from 45 to 170~m below the surface.

The LHC is the final stage of a chain of accelerators that progressively boost proton beam energies before injection, as illustrated in Fig.~\ref{fig:cern_complex}. Protons are produced by ionizing hydrogen gas in Linac4, a linear accelerator that brings them to 160~MeV. They are then passed to the Proton Synchrotron Booster (PSB), which accelerates them to 2~GeV across four superimposed synchrotron rings. The Proton Synchrotron (PS) raises the beam energy to 26~GeV and defines the bunch structure, followed by the Super Proton Synchrotron (SPS) at 450~GeV, which serves as the final injector into the LHC. Inside the LHC, the two counter-rotating beams are accelerated to 6.5~TeV per beam during Run~2 (2015--2018) and 6.8~TeV during Run~3 (2022--2026), following the second long shutdown (LS2, 2019--2021)~\cite{Fartoukh:2790409}. The beams are guided on their circular path by 1232 superconducting dipole magnets operating at 1.9~K with NbTi cables capable of producing fields up to 8.3~T; additional superconducting quadrupole and higher-order magnets provide focusing and orbit correction. The key beam parameters for both running periods are summarized in Table~\ref{tab:lhc_parameters}.

\begin{figure}[htbp]
  \centering
  \includegraphics[width=0.85\textwidth]{Figures/cern_accelerator_complex.pdf}
  \caption{The CERN accelerator complex. Protons are progressively accelerated through Linac4, PSB, PS, and SPS before injection into the LHC~\cite{Mobs:2197559}.}
  \label{fig:cern_complex}
\end{figure}

\begin{table}[htbp]
\centering
\caption{LHC beam parameters during Run~2 and Run~3 operation.}
\label{tab:lhc_parameters}
\begin{tabular}{lcc}
\hline\hline
Parameter & Run~2 & Run~3 \\
\hline
Center-of-mass energy & 13\TeV & 13.6\TeV \\
Beam energy & 6.5\TeV & 6.8\TeV \\
Number of bunches per beam & up to 2556 & up to 2400 \\
Protons per bunch & $\sim 1.15 \times 10^{11}$ & $\sim 1.2 \times 10^{11}$ \\
Bunch spacing & 25~ns & 25~ns \\
Peak instantaneous luminosity [cm$^{-2}$s$^{-1}$] & $2.1 \times 10^{34}$ & $2.0 \times 10^{34}$ \\
\hline\hline
\end{tabular}
\end{table}

The instantaneous luminosity $\mathcal{L}$ characterizes the collision rate per unit cross section:
\begin{equation}
\mathcal{L} = \frac{N_1 N_2 n_b f_{\text{rev}}}{4\pi\sigma_x\sigma_y} F
\end{equation}
where $N_{1,2}$ are the numbers of particles per bunch in each beam, $n_b$ is the number of colliding bunch pairs, $f_{\text{rev}}$ is the revolution frequency, $\sigma_{x,y}$ are the transverse beam sizes at the interaction point, and $F$ is a geometric reduction factor accounting for the crossing angle. The integrated luminosity $L = \int \mathcal{L}(t)\,dt$ quantifies the total collision exposure. During Run~2, CMS recorded a total of 138\fbinv~\cite{CMS:2021xjt,CMS:2018elu,CMS:2019jhp}. Run~3 commenced in 2022 and is ongoing; CMS recorded approximately 300\fbinv through 2025~\cite{CMS:2024onh}, with the full Run~3 dataset projected to reach 500\fbinv by 2026. The cumulative integrated luminosity delivered to CMS is shown in Fig.~\ref{fig:lumi_run2run3}, and the specific dataset used in this analysis is described in Chapter~\ref{ch:datasets}.

The high instantaneous luminosity also produces multiple simultaneous proton-proton interactions per bunch crossing, a phenomenon known as pileup. The mean number of pileup interactions per crossing is:
\begin{equation}
\langle\mu\rangle = \frac{\mathcal{L} \sigma_{\text{inel}}}{n_b f_{\text{rev}}}
\end{equation}
where $\sigma_{\text{inel}} \approx 80$~mb is the inelastic proton-proton cross section. During Run~2, the average pileup ranged from approximately 20 to 50 interactions per bunch crossing, increasing to 40--60 in Run~3. These additional interactions contribute spurious particles and energy deposits that must be mitigated during physics object reconstruction, as discussed in Section~\ref{sec:pf}.

\begin{figure}[htbp]
  \centering
  \includegraphics[width=0.48\textwidth]{Figures/lumi_run2.pdf}
  \includegraphics[width=0.48\textwidth]{Figures/lumi_run3.pdf}
  \caption{Cumulative integrated luminosity recorded by CMS during Run~2 (left)~\cite{CMS:2021xjt} and Run~3 2022--2025 (right)~\cite{CMS:2024onh}.}
  \label{fig:lumi_run2run3}
\end{figure}

\begin{figure}[htbp]
  \centering
  \includegraphics[width=0.48\textwidth]{Figures/inst_lumi_run2.pdf}
  \includegraphics[width=0.48\textwidth]{Figures/inst_lumi_run3.pdf}
  \caption{Peak instantaneous luminosity delivered to CMS per day during Run~2 (left)~\cite{CMS:2021xjt} and Run~3 (right)~\cite{CMS:2024onh}.}
  \label{fig:inst_lumi}
\end{figure}

\begin{figure}[htbp]
  \centering
  \includegraphics[width=0.75\textwidth]{Figures/pileup_run2run3.pdf}
  \caption{Distribution of the mean number of pileup interactions per bunch crossing for each data-taking year in Run~2 and Run~3~\cite{CMS:2021xjt,CMS:2024onh}.}
  \label{fig:pileup}
\end{figure}

\section{The Compact Muon Solenoid Detector}
\label{sec:cms}

The Compact Muon Solenoid (CMS)~\cite{CMS:2008xjf} is a general-purpose particle detector designed to study a wide range of physics processes at the LHC. It is characterized by a powerful superconducting solenoid magnet, fine-grained calorimetry, excellent tracking capabilities, and precise muon detection systems.

\subsection{Detector Overview and Coordinate System}

The CMS detector has a cylindrical geometry 21.6~m in length and 14.6~m in diameter, with a total weight of approximately 14,000 tonnes. Proceeding from the interaction point outward, the main subsystems are the silicon pixel and strip tracker, the electromagnetic calorimeter (ECAL), the hadron calorimeter (HCAL), the superconducting solenoid magnet, and the muon detection systems. A transverse slice of the detector is shown in Fig.~\ref{fig:cms_slice}.

The CMS coordinate system is right-handed with the origin at the nominal interaction point. The $x$-axis points toward the center of the LHC ring, the $y$-axis points vertically upward, and the $z$-axis points along the beam direction toward the Jura mountains. The azimuthal angle $\phi$ is measured from the $x$-axis in the transverse plane, and the polar angle $\theta$ is measured from the $z$-axis. The transverse momentum $\pT = p\sin\theta$ and transverse energy $\ET = E\sin\theta$ are defined perpendicular to the beam axis. The pseudorapidity,
\begin{equation}
\eta = -\ln\left[\tan\left(\frac{\theta}{2}\right)\right],
\end{equation}
is preferred over the polar angle because differences in $\eta$ are invariant under Lorentz boosts along the beam axis. The angular separation between two objects is expressed as $\Delta R = \sqrt{(\Delta\eta)^2 + (\Delta\phi)^2}$, and is used throughout this analysis for jet clustering, lepton isolation, and object matching. The central barrel region covers $|\eta| < 1.5$, while endcap detectors extend coverage to $|\eta| < 2.5$ for tracking and $|\eta| < 5.0$ for calorimetry, as illustrated in Fig.~\ref{fig:cms_eta}.

\begin{figure}[htbp]
  \centering
  \includegraphics[width=0.85\textwidth]{Figures/cms_detector_slice.pdf}
  \caption{Transverse slice of the CMS detector showing the main subsystems and the response of different particle types~\cite{CMS:2008xjf}.}
  \label{fig:cms_slice}
\end{figure}

\begin{figure}[htbp]
  \centering
  \includegraphics[width=0.85\textwidth]{Figures/cms_detector_eta.pdf}
  \caption{Longitudinal cross-section of one quarter of the CMS detector showing the pseudorapidity coverage of each subsystem~\cite{CMS:2008xjf}.}
  \label{fig:cms_eta}
\end{figure}

\subsection{Silicon Tracker}

The CMS silicon tracker~\cite{CMS:2014pgm} is the innermost detector subsystem, designed to measure the trajectories of charged particles emerging from the collision point. It consists of two main components: a high-granularity pixel detector at the innermost radii, optimized for precision vertexing near the interaction point, and a surrounding silicon strip tracker providing large-area coverage, as shown in Fig.~\ref{fig:cms_tracker}.

The pixel detector uses silicon sensors with pixel sizes of $100 \times 150~\mu\text{m}^2$. During Run~2, the original 3-layer system was replaced by a Phase-1 upgraded 4-layer detector in early 2017~\cite{TrackerGroupoftheCMS:2020muy}, comprising four barrel layers (BPIX) at radii of 3.0, 6.8, 10.9, and 16.0~cm, and three forward disks (FPIX) on each side at $|z| = 29.1$, 39.6, and 51.6~cm. The silicon strip tracker surrounds the pixel detector and extends tracking coverage to $|\eta| < 2.5$. It is organized into the Tracker Inner Barrel (TIB, four layers up to $|z| < 65$~cm), the Tracker Inner Disks (TID, three disks per side), the Tracker Outer Barrel (TOB, six layers up to $|z| < 110$~cm), and the Tracker Endcaps (TEC, nine disks per side). The total active silicon area is approximately 200~m$^2$ with about 75 million readout channels. The track reconstruction efficiency for promptly produced charged particles with $\pT > 0.9\GeV$ is approximately 94\% for $|\eta| < 0.9$~\cite{CMS:2014pgm}.

The fine granularity of the pixel detector near the interaction point is essential for reconstructing secondary vertices from long-lived particles. B mesons, for example, have decay lengths of $c\tau \approx 450$--500~$\mu$m; resolving such displaced vertices from the primary interaction point requires a transverse impact parameter resolution well below this scale. For isolated muons in the central region ($|\eta| < 1.4$), the CMS tracker achieves $\sigma(d_0) \approx 10~\mu$m~\cite{CMS:2014pgm}, enabling efficient identification of b-tagged jets important to this analysis. The transverse momentum resolution is $\sigma(\pT)/\pT \approx 1\%$ at $\pT \approx 10\GeV$ and $\approx 2.8\%$ at $\pT \approx 100\GeV$ in the central region, degrading at higher \pT due to increasing curvature measurement uncertainty.

\begin{figure}[htbp]
  \centering
  \includegraphics[width=0.85\textwidth]{Figures/cms_tracker_layout.pdf}
  \caption{Layout of the CMS silicon tracker showing the pixel detector (innermost layers) and the surrounding silicon strip tracker subsystems~\cite{CMS:2014pgm}.}
  \label{fig:cms_tracker}
\end{figure}

\subsection{Calorimeter System}

The CMS calorimeter system consists of two complementary detectors — the electromagnetic calorimeter (ECAL) and the hadron calorimeter (HCAL) — which together measure the energies of electrons, photons, and hadrons produced in pp collisions.

The ECAL~\cite{CMS:2013lxn} uses lead tungstate (PbWO$_4$) crystals, chosen for their short radiation length (0.89~cm), small Molière radius (2.2~cm), and fast scintillation response (80\% of light emitted within 25~ns). The barrel region (EB) houses 61,200 crystals covering $|\eta| < 1.479$ with front-face dimensions of $22 \times 22$~mm$^2$, while the two endcaps (EE) each contain 7,324 crystals covering $1.479 < |\eta| < 3.0$. A preshower detector (ES) consisting of silicon strip sensors is placed in front of the endcap crystals to improve $\pi^0$ rejection. Scintillation light is collected by avalanche photodiodes (APDs) in the barrel and vacuum phototriodes (VPTs) in the endcaps. The energy resolution of the ECAL is parameterized as:
\begin{equation}
\frac{\sigma(E)}{E} = \frac{2.8\%}{\sqrt{E}} \oplus \frac{12\%}{E} \oplus 0.3\%
\end{equation}
where $E$ is in \GeV and $\oplus$ denotes addition in quadrature.

The HCAL~\cite{CMS:2012sof} surrounds the ECAL and measures hadronic shower energies using brass absorber plates interleaved with plastic scintillator tiles. The barrel (HB, $|\eta| < 1.3$) and endcap (HE, $1.3 < |\eta| < 3.0$) sections provide coverage in the central region, supplemented by an outer calorimeter (HO) placed outside the solenoid as a tail catcher, and a forward steel/quartz-fiber calorimeter (HF) extending coverage to $|\eta| < 5.0$. The energy resolution for single hadrons is approximately $\sigma(E)/E \approx 100\%/\sqrt{E} \oplus 5\%$.

During LS2, the HB and HE underwent a significant upgrade for Run~3 operation~\cite{CMS:2023gfb}. The photosensors were replaced from hybrid photodiodes (HPDs) to silicon photomultipliers (SiPMs), offering higher photon detection efficiency, lower noise, and improved radiation tolerance. New front-end readout electronics (QIE11) were installed simultaneously to exploit the finer longitudinal segmentation now available: the HB readout was upgraded from a single depth layer to four, and the HE from three to up to seven layers. This enhanced longitudinal segmentation enables more precise shower-profile reconstruction, improving the separation of electromagnetic and hadronic deposits and the jet energy resolution in Run~3.

\subsection{Superconducting Solenoid}

A strong magnetic field is central to the CMS design, as it curves the trajectories of charged particles and enables momentum measurements from track curvature. The CMS solenoid~\cite{CMS:2009nxa} generates a uniform 3.8~T field inside its bore — among the strongest ever produced in a large detector — with an inner diameter of 6.3~m and a length of 12.5~m. The coil operates at 18.2~kA and stores 2.6~GJ of energy. The magnetic flux is returned through a 12,000-tonne iron yoke, which also serves as the structural support and absorber for the muon detection system described in the following section.

\subsection{Muon System}

The muon system~\cite{CMS:2018rym} is located outside the solenoid and interleaved with the iron return yoke, providing muon identification, triggering, and standalone momentum measurement. Because muons traverse the full detector with minimal energy loss, they serve as clean signatures of many physics processes relevant to this analysis. The system employs three complementary gaseous detector technologies, each suited to the distinct rate and field conditions in different detector regions, as illustrated in Fig.~\ref{fig:cms_muon}.

In the barrel region ($|\eta| < 1.2$), where the muon rate is low and the solenoidal field is uniform, drift tube (DT) chambers are used. Uniform field conditions make electron drift paths predictable and position measurement via drift time reliable. Four DT stations are embedded in the return yoke, achieving a single-hit position resolution of approximately 200~$\mu$m.

The endcap regions ($0.9 < |\eta| < 2.4$) present more demanding conditions: higher muon rates and a non-uniform, non-axial magnetic field that would distort electron drift paths. Cathode strip chambers (CSCs) are employed instead, determining position from the intersection of anode wire and cathode strip signals and thus insensitive to drift path distortions. Each CSC consists of six layers of anode wires interleaved with cathode strips, providing two-dimensional measurements in a single chamber with a resolution of 50--140~$\mu$m depending on the chamber type.

Resistive plate chambers (RPCs) complement both the barrel ($|\eta| < 1.6$) and endcap ($0.9 < |\eta| < 1.9$) regions with a distinct purpose. While DTs and CSCs provide precise spatial information, their timing resolution is insufficient to unambiguously assign a muon to a specific 25~ns-spaced LHC bunch crossing. RPCs sacrifice spatial resolution (${\sim}1$~cm) for excellent timing resolution (${\sim}1$~ns), providing a reliable bunch-crossing tag that anchors the muon trigger decision.

The three technologies are not mutually exclusive: in the overlap region ($0.9 < |\eta| < 1.2$), both DT and CSC chambers are present, providing redundant measurements that improve reconstruction robustness and trigger efficiency.

For Run~3, two upgrades extended the muon system in the forward region~\cite{CMS:2023gfb}. Additional RPC stations (RE3/1 and RE4/1) were installed in the inner endcap disks, extending RPC timing coverage to $|\eta| < 2.1$. Simultaneously, triple-foil gas electron multiplier (GEM) chambers (GE1/1) were installed in the first endcap muon station, covering $1.55 < |\eta| < 2.18$. The forward region experiences the highest particle rates and radiation doses in the muon system; GEM detectors offer superior radiation hardness and rate capability compared to conventional wire chambers, complementing the existing CSCs and improving muon trigger efficiency and fake-rate rejection in this demanding region.

The combined tracker and muon system achieves $\sigma(\pT)/\pT \approx 1$--2\% for $\pT < 100\GeV$ and better than 10\% up to $\pT \approx 1\TeV$.

\begin{figure}[htbp]
  \centering
  \includegraphics[width=0.85\textwidth]{Figures/cms_muon_system.pdf}
  \caption{Longitudinal cross-section of one quarter of the CMS muon system, showing the locations of the DT (barrel), CSC (endcap), RPC (barrel and endcap), and GEM (forward endcap) detector stations~\cite{CMS:2018rym}.}
  \label{fig:cms_muon}
\end{figure}

%% ============================================================
\section{Data Acquisition and Trigger}
\label{sec:daq}
%% ============================================================

At the LHC design luminosity, proton bunches cross at a rate of 40~MHz, producing approximately $10^9$ inelastic interactions per second. Only a small fraction of these collisions contain processes of physics interest, and the full CMS detector readout of roughly 1~MB per event makes it impractical to store every bunch crossing. The CMS data acquisition (DAQ) and trigger system~\cite{CMS:2016ngn,CMS:2020cmk,CMS:2024aqx} selects and retains this small fraction in real time, reducing the 40~MHz input rate to approximately 1~kHz of events written to permanent storage.

The first stage, the Level-1 (L1) trigger~\cite{CMS:2020cmk}, is implemented in custom FPGA-based electronics and must reach a decision within a fixed pipeline latency of approximately 4~$\mu$s, during which the full detector data are buffered on-detector. The L1 system receives coarse-grained trigger primitives from the calorimeters and muon detectors and processes them in two layers of programmable electronics. The calorimeter trigger reconstructs jets, electrons, photons, hadronically decaying $\tau$-leptons, and energy sums ($\HT$, $\MET$), while the muon trigger combines information from the DT, CSC, and RPC systems to form muon candidates. A global trigger (GT) then evaluates a configurable trigger menu — consisting of several hundred algorithms based on object multiplicities, transverse momenta, and their topological combinations — and issues a Level-1 Accept (L1A) for selected events. The L1 trigger reduces the rate from 40~MHz to approximately 100~kHz, the maximum sustainable by the CMS front-end readout electronics.

Events passing the L1 trigger are fully read out by the DAQ system, which assembles data fragments from all detector subsystems via a high-throughput event-builder network at an aggregate throughput of approximately 100~GB/s. These assembled events are passed to the High-Level Trigger (HLT)~\cite{CMS:2024aqx}, a software-based filter running on a farm of commercial multi-core processors. The HLT executes a streamlined version of the offline reconstruction software (CMSSW) organized into trigger paths: sequences of reconstruction and selection modules applied progressively to each event, with early rejection used to minimize processing time. The HLT has full access to the tracker, calorimeter, and muon system data, enabling offline-quality object reconstruction. The HLT reduces the event rate to approximately 1~kHz for permanent storage and offline reconstruction. A data scouting stream additionally stores events at a higher rate (${\sim}30$~kHz) with reduced event content for specific analyses, and a data parking stream retains additional full events for delayed offline reconstruction.

For Run~3, no major hardware upgrade to the L1 trigger was performed, but significant algorithmic improvements were deployed exploiting the finer longitudinal segmentation of the upgraded HCAL and the new GE1/1 GEM chambers~\cite{CMS:2023gfb}. The L1 output rate was increased to approximately 100--110~kHz to accommodate the higher luminosity. The HLT output rate was similarly increased, with approximately 2.6~kHz stored for prompt reconstruction, ${\sim}3$~kHz for parking, and ${\sim}30$~kHz for scouting. A dedicated 40~MHz scouting system, commissioned in the early part of Run~3, further extends the physics reach by recording trigger primitive information at the full bunch-crossing rate. The specific trigger algorithms and \pT thresholds used to select events for this analysis are described in Chapter~\ref{ch:datasets}.

%% ============================================================
\section{Particle-Flow Event Reconstruction}
\label{sec:pf}
%% ============================================================

The CMS detector employs the particle-flow (PF) algorithm~\cite{CMS:2017yfk} to reconstruct and identify every stable particle produced in a collision. Rather than reconstructing physics objects independently from individual subsystems, PF combines information from the tracker, calorimeters, and muon system simultaneously, exploiting the complementary strengths of each detector. The fine-grained silicon tracker, in particular, enables the separation of charged and neutral hadrons in the calorimeters — a distinction that would be impossible from calorimetry alone — and provides precise momentum measurements for charged particles at lower energies than calorimeter clusters.

\subsection*{Element Building and Linking}

The algorithm begins by constructing basic elements independently in each detector subsystem. In the silicon tracker, charged-particle tracks are reconstructed using an iterative algorithm based on Kalman filtering~\cite{Fruhwirth:1987fm}, where each successive iteration progressively relaxes seed requirements to recover tracks with larger impact parameters or lower \pT, improving efficiency for displaced tracks from heavy-flavor decays and soft particles. The 3.8~T solenoidal field bends charged-particle trajectories in the transverse plane, and the sagitta of the resulting helix — measured across the full radial extent of the tracker — directly determines the transverse momentum via $\pT = qBr$, where $r$ is the radius of curvature. Because the sagitta scales as $1/\pT$, the relative momentum resolution $\sigma(\pT)/\pT$ degrades linearly with \pT; in the central region this gives $\sigma(\pT)/\pT \approx 1\%$ at $\pT \approx 10\GeV$ rising to $\approx 2.8\%$ at $\pT \approx 100\GeV$ (see Section~\ref{sec:cms}). This tracker-based momentum measurement is substantially more precise than what the calorimeters alone could provide for charged particles below a few hundred \GeV, and it is the dominant contribution to PF momentum resolution in that regime. In the calorimeters, topological clusters are formed separately in the ECAL and HCAL by seeding at local energy maxima and aggregating neighboring cells above a noise threshold. In the muon system, standalone muon tracks are built from local segments in the DT, CSC, and RPC chambers independently of the tracker.

These elements are then linked into blocks based on geometric proximity. Tracker tracks are extrapolated to the calorimeter surfaces and associated with clusters if the extrapolated position falls within the cluster envelope; ECAL clusters are further linked to HCAL clusters when their positions overlap; and tracker tracks matched to muon system segments form global muon candidates. All elements sharing at least one link are grouped into a block and processed independently in the subsequent identification stage.

\subsection*{Particle Identification}

Within each block, particles are identified in order of decreasing reconstruction purity. The five resulting PF candidate types, and the logic by which each is identified, are as follows.

Muons are identified first, as they provide the most unambiguous detector signature~\cite{CMS:2018rym}. Muon reconstruction proceeds through three successive stages: standalone muon tracks are built from muon system segments alone; tracker muons are formed by extrapolating silicon tracker tracks outward and matching them to at least one muon segment; and global muons are produced by a combined refit of all tracker and muon system hits. A particle-flow muon is then defined by selecting global or tracker muons whose momentum and isolation are consistent with a prompt muon. For the momentum assignment, the tracker measurement is preferred over the global-fit value when the two are in good agreement, as the silicon tracker provides superior resolution for $\pT \lesssim 200\GeV$: the combined tracker-plus-muon-system resolution is $\sigma(\pT)/\pT \approx 1$--$2\%$ in this regime, degrading to below 10\% up to $\pT \approx 1\TeV$ where the muon system lever arm becomes increasingly important (see Section~\ref{sec:cms}). Muon candidates are removed from the block before subsequent steps, and the calorimeter energy deposits they are expected to produce are subtracted.

Electrons and photons are identified next using the ECAL, exploiting the fine granularity and fast response of the lead tungstate crystals~\cite{CMS:2020uim}. Electrons are reconstructed by matching ECAL superclusters — which collect the energy of the primary deposit together with any bremsstrahlung photons emitted along the curved track — to tracks fitted with a Gaussian Sum Filter (GSF) algorithm that explicitly models non-Gaussian energy loss. The electron energy is determined from a combination of the GSF track momentum and the supercluster energy, with the relative weight depending on \pT and the amount of bremsstrahlung. The resulting energy resolution for electrons from $\PZ$ boson decays ranges from 2 to 5\% depending on pseudorapidity and bremsstrahlung fraction~\cite{CMS:2020uim}. Photons, leaving no tracker hit, are identified as compact ECAL clusters not linked to any track; non-isolated photons overlapping with hadronic activity are distinguished from isolated ones by the presence of surrounding HCAL energy.

The remaining calorimeter energy is interpreted as hadrons. Charged hadrons are identified from tracker tracks matched to calorimeter clusters, with any excess cluster energy above the track momentum attributed to overlapping neutral particles. Residual calorimeter deposits with no associated track are classified as neutral hadrons (HCAL-dominated) or non-isolated photons (ECAL-only).

\subsection*{Pileup Mitigation and Higher-Level Reconstruction}

The complete set of PF candidates serves as the input for all higher-level object reconstruction. Jet clustering, lepton isolation calculation, and $\ptvecmiss$ reconstruction are all derived from this unified candidate list, which is one of the key advantages of the PF approach over subsystem-based reconstruction.

Jets are clustered using the anti-$k_{\mathrm{T}}$ algorithm~\cite{Cacciari:2008gp} with a distance parameter of $R = 0.4$, applied to the PF candidate collection after pileup mitigation. The algorithm defines a distance measure between each pair of particles $i$ and $j$ as $d_{ij} = \min(k_{\mathrm{T}i}^{-2}, k_{\mathrm{T}j}^{-2})\,\Delta R_{ij}^2 / R^2$, and between each particle and the beam as $d_{iB} = k_{\mathrm{T}i}^{-2}$. At each step, the smallest distance is found: if it is a $d_{ij}$, the two particles are merged; if it is a $d_{iB}$, the particle is declared a jet and removed. The inverse-$\pT^2$ weighting causes hard particles to cluster their soft neighbors before soft particles cluster each other, so the resulting jet boundaries are determined by the hard cores and are insensitive to soft radiation. This makes the algorithm both infrared- and collinear-safe, and produces jets with a regular, cone-like geometry of radius $R$.

The raw jet energies reconstructed from PF candidates are subject to detector effects and must be corrected before use. Jet energy scale (JES) corrections are applied in sequential layers: the L1 correction removes the average energy contribution from pileup interactions using the FastJet area method; the L2 correction removes $\eta$-dependent non-uniformities in the detector response; and the L3 correction restores the absolute energy scale to the particle level using MC simulation. An additional residual correction derived from data is applied to account for any remaining discrepancy between data and simulation, using dijet and $\gamma$+jet balance methods. After these corrections, the jet energy resolution (JER) in simulation is further smeared to match the resolution measured in data, which is typically 15--20\% at $\pT \approx 30\GeV$ and improves to 5--10\% at $\pT \approx 200\GeV$ in the central region.

Before jet clustering, pileup interactions must be mitigated, as they contribute spurious low-\pT particles that degrade jet energy resolution and $\ptvecmiss$ reconstruction. The primary vertex (PV) is defined as the reconstructed vertex with the largest sum of $\pT^2$ of associated tracks, and serves as the reference point for distinguishing hard-scatter particles from pileup. Charged PF candidates associated with a vertex other than the PV are removed before jet clustering — a procedure known as charged hadron subtraction (CHS). While CHS effectively eliminates the charged component of pileup, it leaves neutral particles unaddressed, as these carry no track and cannot be vertex-associated. For Run~2, jets in this analysis are reconstructed from the CHS-corrected collection (AK4PFchs), with residual pileup jets further suppressed by the pileup jet identification (PU Jet ID) discriminant~\cite{CMS:2020ebo}, a multivariate classifier based on tracking and jet shape variables.

For Run~3, the significantly higher mean pileup of 40--60 interactions per bunch crossing motivates replacing CHS with a more comprehensive approach. The PUPPI algorithm~\cite{Bertolini:2014bba} assigns each PF candidate a weight between 0 and 1 representing the probability that it originates from the hard interaction, computed from the local charged-to-neutral particle activity around each candidate. Because PUPPI weights are applied particle-by-particle before jet clustering, the algorithm simultaneously suppresses both charged and neutral pileup contributions. Jets are reconstructed from the PUPPI-weighted collection (AK4PFPuppi) for all Run~3 data in this analysis.

Jets originating from b quarks are identified by exploiting the long lifetime of b hadrons ($c\tau \approx 450$--500~$\mu$m), which produces charged tracks with significant impact parameters and secondary vertices displaced from the PV, both resolved by the CMS silicon tracker. The CMS b-tagging algorithm has evolved significantly over the course of Run~2. DeepCSV~\cite{CMS:2017wtu} is a fully connected deep neural network with four hidden layers of 100 nodes each, trained on 68 high-level input features derived from the six most displaced tracks and the most displaced secondary vertex in the jet, together with global jet variables. It outputs simultaneous probabilities for b, c, and light-flavor hypotheses. DeepJet~\cite{Bols:2020bkb} supersedes DeepCSV by replacing the fixed set of high-level features with the full low-level information of all jet constituents. Its architecture consists of three parallel LSTM branches — processing charged particles, neutral particles, and secondary vertices separately — whose outputs are concatenated with global jet variables and passed through a series of fully connected layers, yielding approximately 650 input features in total. This design avoids the information loss inherent in preselecting a fixed number of tracks or vertices, and allows the network to learn displacement and topology signatures directly from the constituent level. At a light-flavor jet misidentification rate of 1\%, DeepCSV achieves a b-jet identification efficiency of approximately 68\%~\cite{CMS:2017wtu}. DeepJet provides a significant improvement over this baseline, with relative efficiency gains of up to 20\% at lower mistag rates~\cite{Bols:2020bkb}. This analysis uses DeepJet for b-tagging across both Run~2 and Run~3.

The missing transverse momentum $\ptvecmiss$ is computed as the negative vector sum of the transverse momenta of all PF candidates in the event~\cite{CMS:2019ctu}, and serves as the experimental proxy for undetected particles such as neutrinos. Two variants exist depending on the pileup mitigation applied to the PF input collection. CHS-based $\ptvecmiss$ (PF+CHS MET) uses the charged-hadron-subtracted candidate list, which removes charged pileup contributions but leaves neutral pileup particles intact; these residual neutral deposits introduce a stochastic smearing of the true $\ptvecmiss$ that worsens with increasing pileup. PUPPI-based $\ptvecmiss$ (PUPPI MET) instead uses the PUPPI-weighted candidate list, suppressing both charged and neutral pileup contributions and thus providing better $\ptvecmiss$ resolution under high-pileup conditions. This analysis uses PUPPI MET for both Run~2 and Run~3 datasets: while CHS MET is adequate at the moderate pileup levels of early Run~2, the superior neutral pileup suppression of PUPPI MET provides more consistent resolution across the full pileup range of Run~2 and Run~3, and is the CMS standard for analyses requiring precise $\ptvecmiss$ reconstruction.

The JES corrections described above must also be propagated to $\ptvecmiss$ to maintain global transverse momentum conservation. This is achieved through the type-I correction, which replaces the raw jet momenta in the $\ptvecmiss$ sum with their JES-corrected counterparts, ensuring that any shift in jet energy is consistently reflected in the $\ptvecmiss$ measurement.

Lepton isolation quantifies the hadronic and electromagnetic activity in the vicinity of a lepton candidate, and is the primary handle for distinguishing prompt leptons — produced directly in electroweak decays such as $\PW \to \ell\nu$ or $\PZ \to \ell\ell$ — from nonprompt leptons arising from semileptonic heavy-flavor decays inside jets. In the PF framework, isolation is calculated as the scalar sum of the \pT of all PF candidates within a cone of radius $\Delta R$ around the lepton, divided by the lepton \pT, with an explicit subtraction of the pileup contribution estimated from the local energy density $\rho$. A standard fixed-cone approach with $\Delta R = 0.4$ works well for isolated low-\pT leptons, but becomes inefficient for boosted topologies where the decay products of a heavy resonance are collimated: the fixed cone may overlap with jets from the same hard decay, causing genuine prompt leptons to fail isolation. This analysis instead uses a $\pT$-dependent cone — the mini-isolation variable $I_{\text{mini}}$ — in which the cone shrinks as $10\GeV/\pT$ for leptons above $50\GeV$, keeping the isolation region well-separated from the lepton's own radiation while avoiding overlap with nearby jets. The precise definition and working points are given in Chapter~\ref{ch:objects}.

%% ============================================================
\section{Monte Carlo Simulation}
\label{sec:mc_sim}
%% ============================================================

Monte Carlo (MC) simulation provides the theoretical predictions against which collision data are compared. Simulated event samples are used to model both the signal processes and Standard Model backgrounds, to optimize event selection criteria, and to evaluate systematic uncertainties. This section describes the physics modeling underlying MC event generation, the detector simulation and digitization procedures, and the specific software configurations used for the Run~2 and Run~3 simulation campaigns.

\subsection{Event Generation}
\label{sec:mc_evtgen}

Event generation models the full chain of physical processes occurring in a proton-proton collision. The simulation proceeds through several stages: the hard scattering of partons sampled from the proton structure, parton showering, hadronization, and the modeling of the underlying event. Each stage is described below.

\subsubsection{Parton Distribution Functions}

The cross section for a hard-scattering process in proton-proton collisions is computed using the \textit{factorization theorem}~\cite{Collins:1989gx}, which separates the short-distance partonic interaction from the long-distance structure of the proton:
\begin{equation}
\sigma(pp \to X) = \sum_{a,b} \int dx_a\, dx_b\; f_a(x_a, \mu_F^2)\; f_b(x_b, \mu_F^2)\; \hat{\sigma}_{ab \to X}(\mu_F^2, \mu_R^2),
\label{eq:factorization}
\end{equation}
where $f_{a,b}(x, \mu_F^2)$ are the \textit{parton distribution functions} (PDFs) encoding the probability density of finding parton $a$ (or $b$) carrying a momentum fraction $x$ of the proton at the factorization scale $\mu_F$, and $\hat{\sigma}$ is the partonic cross section computed perturbatively at the renormalization scale $\mu_R$.

PDFs are not calculable from first principles and are instead determined from global fits to experimental data spanning deep inelastic scattering, Drell--Yan, and jet production measurements. Their dependence on the energy scale is governed by the Dokshitzer--Gribov--Lipatov--Altarelli--Parisi (DGLAP) evolution equations~\cite{Dokshitzer:1977sg, Gribov:1972ri, Altarelli:1977zs}, which resum collinear logarithms to all orders. The CMS simulation for both Run~2 and Run~3 uses the NNPDF3.1 set at next-to-next-to-leading order (NNLO) accuracy~\cite{NNPDF:2017mvq}, which employs a neural network parameterization and provides PDF uncertainties through a set of Monte Carlo replicas. The PDFs are accessed at runtime via the \textsc{LHAPDF} library~\cite{Buckley:2014ana}.

\subsubsection{Hard Scattering}

The \textit{hard scattering} is the short-distance interaction between two partons, computed perturbatively as a matrix element (ME) at a fixed order in the strong or electroweak coupling constants, typically at leading order (LO) or next-to-leading order (NLO) accuracy. Two classes of ME generators are widely used in CMS:
\begin{itemize}
\item \textsc{MadGraph5\_aMC@NLO}~\cite{Alwall:2014hca} provides automated computation of tree-level and NLO differential cross sections for arbitrary processes in the Standard Model and beyond. It is the primary generator for multileg processes in CMS.
\item \textsc{Powheg}~\cite{Nason:2004rx, Frixione:2007vw} implements the \textsc{Powheg} method, in which the hardest emission is generated first at NLO accuracy, ensuring a smooth interface with the subsequent parton shower.
\end{itemize}
The output of the ME generator is a set of parton-level events, with each event containing the four-momenta and quantum numbers of the outgoing partons, together with an event weight.

\subsubsection{Parton Shower}

The outgoing partons from the hard scattering undergo \textit{parton showering}, an iterative process that models the successive emission of gluons and quark--antiquark pairs. The parton shower resums the leading soft and collinear logarithms via the DGLAP splitting functions, evolving the parton system from the hard-interaction scale $Q$ down to a cutoff scale of $\mathcal{O}(1\GeV)$, below which perturbation theory ceases to be applicable. Initial-state radiation (ISR) from the incoming partons and final-state radiation (FSR) from the outgoing partons are both modeled.

The two general-purpose parton shower programs used in CMS are \textsc{Pythia}~8~\cite{Sjostrand:2014zea} and \textsc{Herwig}~7~\cite{Bellm:2015jjp}. \textsc{Pythia}~8 employs a $p_\mathrm{T}$-ordered dipole shower, while \textsc{Herwig}~7 uses an angular-ordered shower as its default algorithm. Both programs also handle the subsequent stages of hadronization, underlying event modeling, and particle decays.

\subsubsection{Matching and Merging}

When the ME calculation includes additional radiated partons, care must be taken to avoid double-counting the same phase-space configurations that are also populated by the parton shower. Two principal approaches are employed in CMS:
\begin{itemize}
\item \textit{Matching} combines an NLO ME calculation with the parton shower while preserving the NLO accuracy of inclusive observables. The MC@NLO method~\cite{Frixione:2002ik} and the \textsc{Powheg} method~\cite{Nason:2004rx} are the two standard implementations used in CMS.
\item \textit{Merging} combines tree-level or NLO ME calculations for different parton multiplicities with the parton shower into a single inclusive sample. The MLM scheme~\cite{Alwall:2007fs} is used for LO multileg merging, while the FxFx scheme~\cite{Frederix:2012ps} extends this to NLO accuracy within the \textsc{MadGraph5\_aMC@NLO} framework.
\end{itemize}

\subsubsection{Hadronization and Particle Decays}

At the cutoff scale of the parton shower, the colored partons are converted into color-neutral hadrons through \textit{hadronization}, a nonperturbative process described by phenomenological models. \textsc{Pythia}~8 implements the Lund string model~\cite{Andersson:1983ia}, in which color flux tubes stretched between partons fragment into hadron pairs when the string tension exceeds the pair-production threshold. \textsc{Herwig}~7 employs a cluster model~\cite{Webber:1983if}, in which gluons are split into quark--antiquark pairs and neighboring quarks form color-singlet clusters that decay isotropically into hadrons. Subsequent decays of unstable hadrons, including $\tau$ leptons, $b$ and $c$ hadrons, are handled by the generator itself or by the dedicated \textsc{EvtGen} package~\cite{Lange:2001uf} for heavy-flavor decays.

\subsubsection{Underlying Event}

In addition to the hard scattering, softer interactions between the remaining partons of the colliding protons give rise to the \textit{underlying event} (UE). This includes \textit{multiple parton interactions} (MPI) and the fragmentation of beam remnants. The UE is modeled by \textsc{Pythia}~8 through a $p_\mathrm{T}$-ordered MPI framework, with free parameters determined by fitting to CMS measurements of charged-particle distributions. The CMS~CP5 tune~\cite{CMS:2019csb} is used as the standard configuration for both Run~2 and Run~3 simulation, based on NNPDF3.1 NNLO PDFs with $\alpha_S(M_Z) = 0.118$.

\subsection{Detector Simulation and Digitization}
\label{sec:mc_detsim}

The stable particles produced by the event generator are propagated through a detailed model of the CMS detector implemented in the \textsc{Geant4} toolkit~\cite{GEANT4:2002zbu, Allison:2016lfl}. The simulation models the full detector geometry and material budget, including electromagnetic and hadronic interactions, multiple scattering, photon conversions, nuclear interactions, and energy deposition in active detector elements.

The resulting energy deposits are processed through a \textit{digitization} step that emulates the response of the detector readout electronics, producing simulated raw data in the same format as real collision data. To model the effect of pileup, simulated minimum-bias interactions generated with \textsc{Pythia}~8 are overlaid on each hard-scatter event using a premixing technique, in which pre-digitized pileup events are mixed with the signal event at the digitization stage. The number of overlaid interactions per bunch crossing is sampled from a Poisson distribution whose mean is tuned to match the observed pileup profile in data. After event reconstruction---performed with the same algorithms applied to collision data---the residual differences in the pileup distribution are corrected by reweighting simulated events to match the luminosity-weighted pileup distribution observed in each data-taking period.

\subsection{Run~2 and Run~3 Simulation Configurations}
\label{sec:mc_config}

The detector simulation configurations differ between the two data-taking periods to reflect changes in detector conditions and software. Table~\ref{tab:mc_config} summarizes the key simulation parameters for the Run~2 Ultra Legacy (UL) and Run~3 campaigns.

\begin{table}[htbp]
\centering
\caption{Summary of the Monte Carlo simulation configurations for the Run~2 Ultra Legacy and Run~3 campaigns.}
\label{tab:mc_config}
\begin{tabular}{lll}
\hline
Parameter & Run~2 (Ultra Legacy) & Run~3 (Summer22/23) \\
\hline
CMSSW release    & \texttt{10\_6\_X}  & \texttt{12\_4\_X} / \texttt{13\_X} \\
\textsc{Geant4} version & 10.4--10.6 & 10.7.2 $\to$ 11.1 \\
Geometry description & DDD (native CMS) & DD4hep + VecGeom \\
Physics list     & FTFP\_BERT\_EMM    & FTFP\_BERT\_EMM (optimized) \\
PDF set          & NNPDF3.1 NNLO     & NNPDF3.1 NNLO \\
UE tune          & CP5               & CP5 \\
Pileup method    & Premixing         & Premixing \\
\hline
\end{tabular}
\end{table}

For Run~2, the CMS detector geometry is described using the native DDD (Detector Description Database) framework, while for Run~3 the description has been migrated to the community-developed DD4hep toolkit with the VecGeom geometry navigation library~\cite{Ivanchenko:2021epjconf}. The transition to \textsc{Geant4}~10.7 for Run~3 introduced the FTFP\_BERT\_EMM physics list with optimized tracking in the magnetic field, yielding approximately 25\% reduction in CPU time per event compared to the Run~2 configuration. Further improvements in \textsc{Geant4}~11.1, including a new transportation process with built-in multiple scattering and the G4HepEm sub-library, have been deployed for later Run~3 campaigns~\cite{Ivanchenko:2024epjconf}.

Despite these technical differences, the physics modeling is consistent across the two periods: both use the NNPDF3.1 NNLO PDF set~\cite{NNPDF:2017mvq} and the CMS~CP5 underlying-event tune~\cite{CMS:2019csb}. The specific event generators, cross sections, and MC samples used in this analysis are detailed in Chapter~\ref{ch:datasets}.

%% ============================================================
\section{Offline Software and Computing}
\label{sec:computing}
%% ============================================================

The CMS experiment produces data at rates and volumes that require a globally distributed computing infrastructure for processing, storage, and analysis. This section describes the computing model used by CMS, the software framework and data formats, and the ongoing preparations for the computing challenges of the High-Luminosity LHC.

\subsection{The WLCG and CMS Computing Model}
\label{sec:wlcg}

CMS relies on the Worldwide LHC Computing Grid (WLCG)~\cite{WLCG:2005tdr}, a distributed computing infrastructure comprising over 170 computing centers across more than 40 countries. The CMS computing model~\cite{CMS:2005ctr} organizes these resources into a tiered hierarchy:
\begin{itemize}
\item \textbf{Tier-0} at CERN performs the first-pass (``prompt'') reconstruction of collision data and archives the RAW data to tape. A secondary Tier-0 facility is hosted at the Wigner Research Centre in Budapest.
\item \textbf{Tier-1} sites (seven worldwide) store a complete copy of the RAW and reconstructed data, perform large-scale reprocessing campaigns, and host centrally produced MC simulation samples.
\item \textbf{Tier-2} sites ($\sim$50 worldwide) provide resources for MC production, user analysis, and calibration tasks.
\end{itemize}
Workflows are managed by the WMAgent and GlideinWMS systems, which provide automated job submission and workload distribution across the Grid. Data transfers between sites are coordinated by the PhEDEx (Run~2) and Rucio (Run~3) data management systems. Individual users and groups may also exploit opportunistic resources at high-performance computing (HPC) centers and cloud platforms through the HTCondor workload management system~\cite{Thain:2005condor}.

\subsection{The CMSSW Framework}
\label{sec:cmssw}

All CMS data processing---from online high-level trigger reconstruction through offline analysis---is performed within the CMS Software framework (\textsc{CMSSW})~\cite{Jones:2015cmssw, CMS:2024run3dev}. \textsc{CMSSW} is built on a modular, plugin-based architecture using the \textsc{C++} programming language. It employs an event data model in which each event is a container of data products, and processing is organized as a sequence of modules (producers, filters, and analyzers) configured via \textsc{Python} scripts. The \texttt{cmsDriver.py} tool automates the generation of configuration files for standard production workflows.

Since Run~2, \textsc{CMSSW} supports multi-threaded event processing using Intel Threading Building Blocks (TBB), enabling concurrent processing of multiple events and modules within a single process. This is essential for efficient use of modern multi-core architectures and reduces the memory footprint per job. For Run~3, the framework has been extended to support heterogeneous computing, allowing reconstruction modules to offload computationally intensive tasks to GPUs~\cite{Bocci:2020hetero}.

\subsection{Data Formats and Processing Chain}
\label{sec:datatiers}

CMS defines a hierarchy of data formats (``data tiers'') of decreasing size and increasing abstraction, designed to balance storage cost against the information content needed at each stage of the analysis workflow:
\begin{itemize}
\item \textbf{RAW}: The unprocessed detector readout after online formatting and trigger decisions. A typical RAW event occupies $\sim$1--2~MB.
\item \textbf{RECO}: The output of full event reconstruction, containing all reconstructed physics objects together with the hits and clusters from which they are built. A RECO event is $\sim$2--4~MB.
\item \textbf{AOD} (Analysis Object Data): A reduced subset of RECO retaining the physics objects and associated metadata needed for most analyses ($\sim$400~kB/event).
\item \textbf{MiniAOD}~\cite{Petrucciani:2015gjw}: Introduced for Run~2, this format stores high-level physics objects with sufficient information for the majority of analyses at $\sim$40~kB/event, a factor of ten reduction relative to AOD.
\item \textbf{NanoAOD}~\cite{Rizzi:2019epjconf}: A flat \textsc{ROOT}~\cite{Brun:1997pa} ntuple format containing only scalar and vector branches for reconstructed objects, at $\sim$1--2~kB/event. NanoAOD can be analyzed without the \textsc{CMSSW} framework, enabling the use of lightweight analysis tools.
\end{itemize}
For simulated samples, the corresponding tiers are labeled with a ``SIM'' suffix (e.g., MINIAODSIM, NANOAODSIM). The production chain for MC samples proceeds as:
\begin{equation*}
\text{GEN-SIM} \to \text{DIGI-RAW} \to \text{RECO} \to \text{MINIAODSIM} \to \text{NANOAODSIM},
\end{equation*}
where GEN-SIM combines event generation with \textsc{Geant4} detector simulation, and DIGI-RAW performs digitization with pileup overlay. Each step is executed as a separate campaign managed through the Monte Carlo Management (McM) system. For collision data, the corresponding chain is:
\begin{equation*}
\text{RAW} \to \text{RECO} \to \text{AOD} \to \text{MiniAOD} \to \text{NanoAOD}.
\end{equation*}

\subsection{Computing for the High-Luminosity LHC}
\label{sec:hllhc_computing}

The High-Luminosity LHC (HL-LHC), expected to begin operations in 2030, will deliver an instantaneous luminosity of up to $7.5 \times 10^{34}\percms$ with an average of approximately 200 simultaneous pileup interactions per bunch crossing~\cite{CMS:2025phase2}. This will increase both the event complexity and the data volume by roughly an order of magnitude relative to Run~3. Current projections indicate that the computing resources required by CMS under the HL-LHC conditions would exceed the expected budget by a factor of 10--20 if no improvements to software efficiency are made~\cite{CMS:2022note008, Albrecht:2018hsf}. Addressing this challenge requires concurrent advances along multiple fronts.

\subsubsection{Heterogeneous Computing}

A central element of the CMS computing strategy for the HL-LHC is the exploitation of \textit{heterogeneous architectures}---in particular, Graphics Processing Units (GPUs) and other accelerators---to supplement traditional CPU-based processing.

\paragraph{GPU-accelerated reconstruction.}
The Patatrack project has demonstrated that pixel track and vertex reconstruction can be efficiently offloaded to GPUs using a cellular automaton algorithm, achieving throughput equivalent to two full CPU nodes on a single NVIDIA T4 GPU~\cite{Bocci:2020patatrack}. Building on this, the CMS high-level trigger farm for Run~3 has been equipped with GPU coprocessors, enabling GPU-based pixel reconstruction, ECAL reconstruction, and HCAL reconstruction to run in production~\cite{CMS:2024run3dev}. The SONIC (Services for Optimized Network Inference on Coprocessors) framework~\cite{CMS:2024sonic} further extends heterogeneous computing by offloading machine-learning inference to remote GPU servers, decoupling the accelerator hardware from the processing nodes.

\paragraph{GPU-accelerated event generation.}
The Madgraph4GPU project~\cite{Valassi:2025cudacpp, Hageboeck:2023mg4gpu} ports the matrix-element computation engine of \textsc{MadGraph5\_aMC@NLO} to GPUs and vectorized CPUs via the CUDACPP plugin. The CMS integration of this tool~\cite{Choi:2025mg4gpu} has demonstrated speedups of up to 10$\times$ on GPUs and 3$\times$ on vectorized CPUs for gridpack production, and up to 7$\times$ (GPU) and 1.5$\times$ (vectorized CPU) for event generation, benchmarked on Drell--Yan and top quark pair production processes. This work represents one of the first demonstrations of GPU-accelerated Monte Carlo event generation within a full LHC experiment production workflow, and is expected to yield significant resource savings as production campaigns scale toward the HL-LHC era.

\subsubsection{Machine-Learning-Accelerated Reconstruction}

Machine learning (ML) techniques offer the potential to improve both the speed and quality of event reconstruction. The CMS \textit{Machine-Learned Particle-Flow} (MLPF) algorithm~\cite{CMS:2026mlpf} replaces the traditional rule-based particle-flow reconstruction with a unified graph neural network model, demonstrating improved jet energy resolution while running efficiently on GPUs. For the Phase-2 high-granularity calorimeter endcap (HGCAL), ML-based clustering and energy regression algorithms are being developed to handle the significantly increased channel count and event complexity~\cite{CMS:2025phase2}.

\subsubsection{Data Management and Analysis Facilities}

The WLCG storage model is evolving from strict site-based data placement toward a more flexible \textit{data lake} architecture~\cite{Bird:2019datalake}, in which large storage facilities serve data to smaller computing sites through caching layers, reducing the total number of replicas while maintaining data accessibility. The shift toward the compact NanoAOD format has enabled a new paradigm of \textit{analysis facilities}---centralized or federated computing platforms that provide interactive, low-latency access to analysis-ready datasets. These facilities, combined with columnar analysis frameworks and scalable task-based schedulers, aim to reduce the turnaround time for physics analyses from days to minutes, a capability that will be essential for the HL-LHC physics program.
